
\chapter{Burgers equation}
The Burgers equation reads (\citet{Debnath2008}):
\begin{align}
\boxed{
\pdiff{u}{t} + u \pdiff{u}{x} - \nu \pdiff[2]{u}{x}= 0, \quad u >0
}
\end{align}
with $\nu$ constant in time and space.
The Burgers equation is written in conservative form, because we will use a finite volume method to solve this equation, it reads:
\begin{align}
    \pdiff{u}{t} + \half \pdiff{u^2}{x} - \pdiff{}{x} \left( \nu \pdiff{u}{x} \right) = 0
\end{align}
Applying the finite volume approach:
\begin{align}
    \int_\Omega \pdiff{u}{t}\, d\Omega + \half \int_\Omega \pdiff{u^2}{x}\, d\Omega \ - \int_\Omega  \pdiff{}{x} \left( \nu \pdiff{u}{x} \right)\, d\Omega = 0
\end{align}
and after applying Green's theorem it yields:
\begin{align}
    \int_\Omega \pdiff{u}{t}\, d\Omega + \half \oint_\Gamma u^2\, d\Gamma
    - \oint_\Gamma  \nu \pdiff{u}{x} \, d\Gamma = 0
\end{align}

%------------------------------------------------------------------------------
\section{Discretization in time and space}
The discretization in time and the discretization in space are discussed in separate sections.
Due to the applying regularization procedure the viscosity ($\nu$) is increased by an artificial viscosity ($\Psi$), so the the implementation of $\nu + \Psi$ is a function of $x$.
%------------------------------------------------------------------------------
\subsection{Discretization in time}
The used discretization in time reads:
\begin{align}
    \int_\Omega \pdiff{u}{t}\, d\Omega & \approx
    \Dx  \Dt_{inv}  \left(\frac{1}{8} \Delta u^{n+1,p+1}_{i-1} + \frac{6}{8} \Delta u^{n+1,p+1}_{i} + \frac{1}{8} \Delta u^{n+1,p+1}_{i+1}\right) +
    \nonumber \\*
    &+ \Dx \Dt_{inv} \left( \frac{1}{8}\left( u^{n+1,p}_{i-1} - u^{n}_{i-1} \right) + \frac{6}{8}\left( u^{n+1,p}_{i} - u^{n}_{i} \right) + \frac{1}{8}\left( u^{n+1,p}_{i+1} - u^{n}_{i+1} \right) \right)
\end{align}
\textbf{Special attention is needed if $\nu$ is time dependent!}
\todo{What to do with time dependent viscosity}
\newpage
%------------------------------------------------------------------------------
\subsection{Discretization in space}
%------------------------------------------------------------------------------
\subsubsection*{Convection-term}
\begin{align}
   \half \oint_\Gamma u^2\, d\Gamma & =
   \half \left. u^2\right|_{i+\half}
   - \half \left. u^2\right|_{i-\half}
\end{align}
with the approximation  $u^2 = (u^{n+\theta, p})^2
+ 2 u^{n+\theta, p} \theta \Delta u$ the discretization reads:
\begin{align}
    & \approx
    u^{n+\theta, p}_{i+\half} \theta \Delta u_{i+\half}
     - u^{n+\theta, p}_{i-\half} \theta \Delta u_{i-\half}
    + \half \left(u^{n+\theta, p}_{i+\half}\right)^2
    - \half \left(u^{n+\theta, p}_{i-\half}\right)^2
\end{align}
%------------------------------------------------------------------------------
\subsubsection*{Artificial viscosity term $\bm \Psi$}
Discretization of $\Psi(x,t)$ in time and space needed

\notyet
%------------------------------------------------------------------------------
\subsubsection*{Viscosity term $\bm \nu + \Psi$}
With $\eps = \nu + \Psi$ the viscosity term reads:
\begin{align}
     - \oint_\Gamma  \eps \pdiff{u}{x} \, d\Gamma =
     - \left(
     \left(\eps \pdiff{u}{x}\right)_{i+\half}
     - \left(\eps \pdiff{u}{x}\right)_{i-\half}
     \right)
\end{align}
approximated by
\begin{align}
     & \Rightarrow
    -\theta \frac{\eps_{i+\half}}{\Dx} \Delta u^{n+1, p+1}_{i+1}
    + \theta \left( \frac{\eps_{i+\half}}{\Dx}
    + \frac{\eps_{i-\half}}{\Dx} \right) \Delta u^{n+1, p+1}_{i}
    - \theta \frac{\eps_{i-\half}}{\Dx} \Delta u^{n+1, p+1}_{i-1}
    +
    \\*
    &\qquad -\left\{ \left( \eps_{i+\half} \frac{u^{n+\theta,p}_{i+1} - u^{n+\theta,p}_{i}}{\Dx} - \eps_{i-\half} \frac{u^{n+\theta,p}_{i} - c^{n+\theta,p}_{i-1}}{\Dx}\right) \right\}
\end{align}
%------------------------------------------------------------------------------
\subsubsection*{Artificial viscosity term}
\notyet

%%------------------------------------------------------------------------------
%\section{Energy Equation}
%
%The \textbf{inviscid} Burgers' equation is given by:
%\begin{align}
%\frac{\partial u}{\partial t} + u \frac{\partial u}{\partial x} = 0.
%\end{align}
%The kinetic energy per unit mass is:
%\begin{align}
%    E = \frac{1}{2} u^2
%\end{align}
%The time derivative of the energy reads:
%\begin{align}
%    \frac{\partial E}{\partial t} = \frac{\partial}{\partial t} \left( \frac{1}{2} u^2 \right) = u \frac{\partial u}{\partial t} =
%    u \left(- u \frac{\partial u}{\partial x} \right) = - u^2 \frac{\partial u}{\partial x} =
%    - \frac{\partial}{\partial x} \left( \frac{1}{3} u^3 \right)
%\end{align}
%So the energy equation becomes:
%\begin{align}
%    \frac{\partial E}{\partial t} + \frac{\partial}{\partial x} \left( \frac{1}{3} u^3 \right) = 0
%\end{align}
%%The error-function reads:
%%\begin{align}
%%    {\mathit Err}\left( \sqrt[3]{\frac{1}{3}\uhat^3} \right) & =  \sqrt[3]{\frac{1}{3}}\left( \uhat - \ubar \right) =
%%    \\*
%%    & = \sqrt[3]{\frac{1}{3}}\left( \ubar + \frac{1}{8}\Dx^2 \pdiff[2]{\ubar}{x} - \ubar \right) =
%%    \\*
%%    & = \sqrt[3]{\frac{1}{3}}\left(  \frac{1}{8}\Dx^2 \pdiff[2]{\ubar}{x} \right) =
%%\end{align}
%This shows that energy is conserved in the inviscid Burgers equation.
%------------------------------------------------------------------------------
\section{Boundary}
Because it is a one dimensional equation we have only boundary conditions at the west- and east-side of the domain.
We consider only a flow direction in positive $x$-direction, so $u>0$.

%------------------------------------------------------------------------------
\subsection{West boundary}
This is an inflow boundary ($u>0$).
\begin{align}
    u^{n+1} &= f(t)
    \\
    u^{n+\theta, p} + \theta \Delta u  &= f(t)
    \\
    \theta \Delta u & = f(t) - u^{n+\theta, p}
\end{align}

%------------------------------------------------------------------------------
\subsection{East boundary}
This is an outflow boundary ($u>0$).

\notyet
%================================================================================

